\scanpage{12}
\mspage{9}
\deleted{\unsure{Prop.}} Notation pour $x\in L$
\[
[x_0,x]=\left\{
\begin{array}{ll}
\{x_0\}&\text{si }x=x_0\\
L=\tilde I&\text{si }x=x_1\\
\tilde I_{x_0,x}&\text{si }x\in\tilde I^\circ=L\setminus\{x_0,x_1\}
\end{array}
\right.
\]
On définit de $=$
\[
[x,x_1]=\left\{\ldots\right.
\]
(en échangeant les rôles de $x_0,x_1$)

\uline{Proposition} Soient $x,y\in L$. Conditions équivalentes

(i) $[x_0,x]\subset[x_0,y]$

(ii) \deleted{$\unsure{\exists}$} $x\in[x_0,y]$

(i bis) $[x,x_1]\supset[y,x_1]$

(ii bis) \deleted{$\unsure{x\in}$} $y\in[x,x_1]$

Si on \deleted{\unsure{app}} désigne par
\[
x\leq y
\]
cette relation en $x,y$, on trouve une relation de \deleted{\unsure{un}} préordre [évident sur (i) ou (i bis)] et $=$ d'ordre
(puisque $[x_0,x]=[x_0,y]$
\deleted{$\unsure{\Longrightarrow\tilde I_{x_0,x}=I_{x_0,y}}$}
\inserted{ou bien $x=x_0,y=x_0$}
\deleted{d'où $x=y$}
si les deux \deleted{$\unsure{\Longrightarrow}$} sont réduits à $x_0$, ou bien $I_{x_0,x}=I_{x_0,y}$ d'où en passant aux $\delta$ $x=y$)
\mspage{11}

\rotatebox{90}{\reflectbox{\unsure{EXTRA STRONG}}}
\rotatebox{90}{\reflectbox{\unsure{EXTRA STRONG}}}
